% Capitolo 2: Fondamenti teorici della polarimetria

\section{Polarizzazione della luce}
\label{sec:polarizzazione}

La luce è un'onda elettromagnetica trasversale. Per un'onda monocromatica che si propaga lungo l'asse $z$, il campo elettrico può essere decomposto in due componenti ortogonali $E_x$ e $E_y$, ciascuna con ampiezza e fase proprie~\cite{hecht2017optics}. La differenza di fase $\delta = \delta_y - \delta_x$ tra le due componenti determina lo stato di polarizzazione: quando $\delta = 0$ o $\pi$ la polarizzazione è lineare, quando $\delta = \pm\pi/2$ con ampiezze uguali è circolare, e nel caso generale è ellittica. Le sorgenti comuni emettono luce non polarizzata, in cui $\delta$ varia casualmente nel tempo; la sovrapposizione di una componente polarizzata e una non polarizzata produce luce parzialmente polarizzata, situazione tipica in ambito sperimentale.


\section{Formalismo di Stokes}
\label{sec:stokes}

Per descrivere la polarizzazione della luce esistono diversi formalismi. Il \emph{formalismo di Jones}, basato su un vettore complesso a due componenti, è efficace ma limitato alla luce completamente polarizzata. Per trattare anche la luce parzialmente polarizzata, situazione comune nella pratica sperimentale, è necessario un formalismo più generale: il \textbf{formalismo di Stokes}, introdotto da George Gabriel Stokes nel 1852~\cite{goldstein2011polarized}.

\subsection{Definizione dei parametri di Stokes}
\label{subsec:stokes_def}

I quattro parametri di Stokes si definiscono a partire dalle componenti del campo elettrico attraverso medie temporali~\cite{collett2005polarized}:
\begin{equation}
\label{eq:stokes_campi}
\Svec = \begin{pmatrix} S_0 \\ S_1 \\ S_2 \\ S_3 \end{pmatrix} = \begin{pmatrix} \langle E_{0x}^2 \rangle + \langle E_{0y}^2 \rangle \\ \langle E_{0x}^2 \rangle - \langle E_{0y}^2 \rangle \\ 2\langle E_{0x} E_{0y} \cos\delta \rangle \\ 2\langle E_{0x} E_{0y} \sin\delta \rangle \end{pmatrix},
\end{equation}
dove le parentesi angolari indicano la media temporale. $S_0$ rappresenta l'intensità totale del fascio; $S_1$ quantifica la prevalenza della polarizzazione orizzontale ($S_1 > 0$) o verticale ($S_1 < 0$); $S_2$ la prevalenza della polarizzazione a $+45^\circ$ ($S_2 > 0$) o $-45^\circ$ ($S_2 < 0$); $S_3$ la prevalenza della polarizzazione circolare destra ($S_3 > 0$) o sinistra ($S_3 < 0$). I parametri di Stokes sono quantità misurabili, proporzionali alle intensità registrate da combinazioni appropriate di polarizzatori e lamine:
\begin{equation}
\label{eq:stokes_intensita}
\Svec = \begin{pmatrix} I_H + I_V \\ I_H - I_V \\ I_{+45} - I_{-45} \\ I_{RCP} - I_{LCP} \end{pmatrix},
\end{equation}
dove $I_H$, $I_V$, $I_{+45}$, $I_{-45}$ sono le intensità trasmesse da polarizzatori lineari orientati rispettivamente a $0^\circ$, $90^\circ$, $+45^\circ$ e $-45^\circ$, mentre $I_{RCP}$ e $I_{LCP}$ corrispondono alle componenti circolare destra e sinistra.

Per qualsiasi stato di polarizzazione vale la disuguaglianza fondamentale:
\begin{equation}
\label{eq:stokes_disuguaglianza}
S_0^2 \geq S_1^2 + S_2^2 + S_3^2,
\end{equation}
con l'uguaglianza valida esclusivamente per luce totalmente polarizzata. Il \textbf{grado di polarizzazione} (DOP, dall'inglese \textit{Degree of Polarization}) quantifica la frazione di luce polarizzata~\cite{bornwolf1999}:
\begin{equation}
\label{eq:DOP}
\text{DOP} = \frac{\sqrt{S_1^2 + S_2^2 + S_3^2}}{S_0},
\end{equation}
e vale 1 per luce totalmente polarizzata e 0 per luce completamente non polarizzata.

Dal punto di vista sperimentale, $S_0$, $S_1$ e $S_2$ si ottengono da misure di intensità attraverso un analizzatore lineare a diversi angoli, come mostrato dalla legge di Malus generalizzata (\cref{sec:malus}). Il parametro $S_3$ richiede l'inserimento di una lamina $\lambda/4$ prima dell'analizzatore.

\subsection{La sfera di Poincaré}
\label{subsec:poincare}

I parametri di Stokes normalizzati $s_1 = S_1/S_0$, $s_2 = S_2/S_0$, $s_3 = S_3/S_0$ definiscono un punto nello spazio tridimensionale. Gli stati totalmente polarizzati giacciono sulla superficie di una sfera unitaria --- la \textbf{sfera di Poincaré}~\cite{goldstein2011polarized} --- il cui equatore rappresenta la polarizzazione lineare e i cui poli rappresentano la polarizzazione circolare destra e sinistra. La distanza dal centro è pari al grado di polarizzazione. Questa rappresentazione geometrica sarà utile nell'interpretazione delle mappe di Stokes del \cref{ch:risultati}.


\section{Legge di Malus generalizzata}
\label{sec:malus}

La legge di Malus classica descrive l'intensità trasmessa da un polarizzatore ideale quando luce linearmente polarizzata lo attraversa. Se l'angolo tra la direzione di polarizzazione della luce incidente e l'asse di trasmissione dell'analizzatore è $\theta$, l'intensità trasmessa vale~\cite{hecht2017optics}:
\begin{equation}
\label{eq:malus_classica}
I(\theta) = I_0 \cos^2\theta.
\end{equation}

Nel caso di luce parzialmente polarizzata, la legge si generalizza utilizzando il formalismo di Stokes. L'intensità trasmessa da un analizzatore lineare ideale orientato a un angolo $\theta$ è~\cite{goldstein2011polarized}:
\begin{equation}
\label{eq:malus_generalizzata}
I(\theta) = \frac{1}{2}\bigl(S_0 + S_1\cos 2\theta + S_2\sin 2\theta\bigr).
\end{equation}
L'\cref{eq:malus_generalizzata} può essere derivata applicando la matrice di Mueller di un polarizzatore lineare ruotato di $\theta$ al vettore di Stokes incidente e leggendo la componente $S_0'$ del vettore uscente (il formalismo delle matrici di Mueller è introdotto nella \cref{sec:mueller}). Si osservi che l'intensità trasmessa dipende solo dai primi tre parametri di Stokes e non da $S_3$: un polarizzatore lineare è insensibile alla polarizzazione circolare.

Acquisendo $N$ immagini a diversi angoli $\theta_i$ dell'analizzatore, l'\cref{eq:malus_generalizzata} fornisce un sistema sovradeterminato di equazioni lineari nei tre parametri incogniti $S_0$, $S_1$ e $S_2$, risolvibile tramite il metodo della pseudo-inversa.


\section{Matrici di Mueller}
\label{sec:mueller}

Il formalismo delle matrici di Mueller descrive la trasformazione dello stato di polarizzazione della luce al passaggio attraverso un elemento ottico. Il vettore di Stokes uscente $\Svec_\text{out}$ è legato a quello incidente $\Svec_\text{in}$ dalla relazione~\cite{goldstein2011polarized}:
\begin{equation}
\label{eq:mueller}
\Svec_\text{out} = \Mmat \, \Svec_\text{in},
\end{equation}
dove $\Mmat$ è una matrice reale $4 \times 4$ caratteristica dell'elemento ottico. Per un sistema composto da $N$ elementi ottici in sequenza, la matrice complessiva si ottiene per moltiplicazione:
\begin{equation}
\label{eq:mueller_composizione}
\Mmat_\text{tot} = \Mmat_N \cdots \Mmat_2 \, \Mmat_1,
\end{equation}
dove $\Mmat_1$ è la matrice del primo elemento incontrato dal fascio.

\subsection{Polarizzatore lineare ideale}
\label{subsec:mueller_pol}

La matrice di Mueller di un polarizzatore lineare ideale con asse di trasmissione orizzontale è:
\begin{equation}
\label{eq:mueller_pol}
\Mmat_\text{pol}(0^\circ) = \frac{1}{2}\begin{pmatrix} 1 & 1 & 0 & 0 \\ 1 & 1 & 0 & 0 \\ 0 & 0 & 0 & 0 \\ 0 & 0 & 0 & 0 \end{pmatrix}.
\end{equation}
Per un polarizzatore con asse ruotato di un angolo $\theta$, la matrice si ottiene applicando la trasformazione di rotazione $\Mmat_\text{pol}(\theta) = \Mmat_R(-\theta)\,\Mmat_\text{pol}(0^\circ)\,\Mmat_R(\theta)$, dove $\Mmat_R(\theta)$ è la matrice di rotazione nello spazio di Stokes (\cref{eq:mueller_rotatore}).

\subsection{Lamina a ritardo (retarder)}
\label{subsec:mueller_retarder}

Una lamina a ritardo introduce una differenza di fase $\delta$ tra le componenti del campo polarizzate lungo l'asse veloce e l'asse lento. La matrice di Mueller per una lamina con asse veloce orizzontale è~\cite{collett2005polarized}:
\begin{equation}
\label{eq:mueller_retarder}
\Mmat_\text{ret}(\delta) = \begin{pmatrix} 1 & 0 & 0 & 0 \\ 0 & 1 & 0 & 0 \\ 0 & 0 & \cos\delta & \sin\delta \\ 0 & 0 & -\sin\delta & \cos\delta \end{pmatrix}.
\end{equation}
I casi di interesse principale sono la lamina \textbf{quarto d'onda} ($\delta = \pi/2$) e la lamina \textbf{mezzo d'onda} ($\delta = \pi$). Per una lamina con asse veloce ruotato di un angolo $\alpha$, si applica la stessa trasformazione di rotazione descritta per il polarizzatore.

\subsection{Rotatore ottico}
\label{subsec:mueller_rotatore}

Un rotatore ottico ruota il piano di polarizzazione di un angolo $\phi$ senza modificare l'ellitticità. La sua matrice di Mueller è~\cite{goldstein2011polarized}:
\begin{equation}
\label{eq:mueller_rotatore}
\Mmat_R(\phi) = \begin{pmatrix} 1 & 0 & 0 & 0 \\ 0 & \cos 2\phi & \sin 2\phi & 0 \\ 0 & -\sin 2\phi & \cos 2\phi & 0 \\ 0 & 0 & 0 & 1 \end{pmatrix}.
\end{equation}
Il fattore $2\phi$ riflette la periodicità $\pi$ della polarizzazione lineare: una rotazione fisica di $\phi$ corrisponde a uno spostamento di $2\phi$ sulla sfera di Poincaré.


\section{Fenomeni che influenzano la polarizzazione}
\label{sec:fenomeni}

I campioni analizzati in questa tesi interagiscono con la polarizzazione attraverso tre meccanismi distinti, ciascuno descritto da una specifica matrice di Mueller.

\textbf{Birifrangenza.} Un materiale è birifrangente quando il suo indice di rifrazione dipende dalla direzione di polarizzazione. Nei cristalli uniassici si distinguono l'indice \textbf{ordinario} $n_o$ e l'indice \textbf{straordinario} $n_e$; la differenza $\Delta n = n_e - n_o$ è la birifrangenza del materiale~\cite{hecht2017optics}. Quando un'onda polarizzata attraversa un mezzo birifrangente di spessore $d$, le due componenti del campo accumulano una differenza di fase detta \textbf{ritardanza}:
\begin{equation}
\label{eq:ritardanza}
\delta = \frac{2\pi}{\lambda}\,\Delta n\, d,
\end{equation}
dove $\lambda$ è la lunghezza d'onda nel vuoto. Le \textbf{lamine d'onda} sfruttano questo effetto: una lamina $\lambda/4$ ($\delta = \pi/2$) converte polarizzazione lineare a $45^\circ$ in circolare, una lamina $\lambda/2$ ($\delta = \pi$) ruota il piano di polarizzazione di un angolo doppio rispetto all'orientazione della lamina. Poiché $\delta$ dipende inversamente da $\lambda$, una lamina calibrata a una data lunghezza d'onda introduce ritardi diversi alle altre lunghezze d'onda --- un aspetto che avrà conseguenze dirette sulla misura di $S_3$ con il sensore a tre canali colore (\cref{ch:analisi}). Una lamina \textbf{zero-order}, il cui ritardo è ottenuto con un singolo passaggio attraverso un cristallo sottile, presenta una dipendenza dalla lunghezza d'onda più debole rispetto a una lamina \textbf{multi-order}.

\textbf{Attività ottica.} Alcuni materiali chirali ruotano il piano di polarizzazione senza alterare l'ellitticità: nel formalismo di Mueller, l'effetto corrisponde a un rotatore (\cref{eq:mueller_rotatore}), con $S_1$ e $S_2$ che ruotano mentre $S_3$ resta invariato. Per soluzioni di sostanze otticamente attive, l'angolo di rotazione segue la \textbf{legge di Biot}~\cite{hecht2017optics}:
\begin{equation}
\label{eq:biot}
\alpha_\text{rot} = [\alpha]_T^\lambda \, l \, c,
\end{equation}
dove $[\alpha]_T^\lambda$ è la \textbf{rotazione specifica} (dipendente dalla temperatura $T$ e dalla lunghezza d'onda $\lambda$), $l$ è il cammino ottico in decimetri e $c$ la concentrazione in \si{\gram\per\milli\liter}. Il campione di saccarosio del \cref{ch:campioni} è stato scelto per verificare questa proprietà e la sua dispersione rotatoria, che in prima approssimazione segue il modello di Drude $[\alpha]_T^\lambda \propto 1/\lambda^2$.

\textbf{Fotoelasticità.} Un materiale trasparente e normalmente isotropo diventa birifrangente quando sottoposto a sollecitazione meccanica~\cite{goldstein2011polarized}. La birifrangenza indotta è proporzionale alla differenza delle tensioni principali attraverso la \textbf{legge tensione-ottica}:
\begin{equation}
\label{eq:fotoelastica}
\Delta n = C\,(\sigma_1 - \sigma_2),
\end{equation}
dove $C$ è la costante fotoelastica del materiale e $\sigma_1$, $\sigma_2$ sono le tensioni principali nel piano perpendicolare alla direzione di propagazione. La ritardanza del campione diventa quindi:
\begin{equation}
\label{eq:ritardanza_fotoelastica}
\delta = \frac{2\pi}{\lambda}\,C\,(\sigma_1 - \sigma_2)\,d.
\end{equation}
Quando un campione fotoelastico è osservato in un polarimetro a campo pieno, si osservano \textbf{frange isocromatiche} (loci a ritardanza costante) e \textbf{isocline} (loci in cui le direzioni principali sono allineate agli assi del polarimetro). Le mappe di ritardanza prodotte dal polarimetro a immagine rendono visibili entrambe le famiglie di frange, come si vedrà nel \cref{ch:risultati}.
